\documentclass[10pt,a4paper]{amsart}
\usepackage{amsmath,amssymb,amsthm,graphicx}
\usepackage{verbatim}
\newenvironment{reference}{\comment}{\endcomment}
%\newenvironment{reference}{}{}
\newenvironment{slogan}{\comment}{\endcomment}
\newenvironment{history}{\comment}{\endcomment}

% For commutative diagrams we use Xy-pic
\usepackage[all]{xy}

% We use 2cell for 2-commutative diagrams.
\xyoption{2cell}
% Theorem environments.
%
\theoremstyle{plain}
\newtheorem{theorem}[subsection]{Theorem}
\newtheorem{proposition}[subsection]{Proposition}
\newtheorem{lemma}[subsection]{Lemma}

\theoremstyle{definition}
\newtheorem{definition}[subsection]{Definition}
\newtheorem{example}[subsection]{Example}
\newtheorem{exercise}[subsection]{Exercise}
\newtheorem{situation}[subsection]{Situation}

\theoremstyle{remark}
\newtheorem{remark}[subsection]{Remark}
\newtheorem{remarks}[subsection]{Remarks}

\numberwithin{equation}{subsection}

% Macros
%
\def\lim{\mathop{\mathrm{lim}}\nolimits}
\def\colim{\mathop{\mathrm{colim}}\nolimits}
\def\Spec{\mathop{\mathrm{Spec}}}
\def\Hom{\mathop{\mathrm{Hom}}\nolimits}
\def\Ext{\mathop{\mathrm{Ext}}\nolimits}
\def\SheafHom{\mathop{\mathcal{H}\!\mathit{om}}\nolimits}
\def\SheafExt{\mathop{\mathcal{E}\!\mathit{xt}}\nolimits}
\def\Sch{\mathit{Sch}}
\def\Mor{\mathop{\mathrm{Mor}}\nolimits}
\def\Ob{\mathop{\mathrm{Ob}}\nolimits}
\def\Sh{\mathop{\mathit{Sh}}\nolimits}
\def\NL{\mathop{N\!L}\nolimits}
\def\CH{\mathop{\mathrm{CH}}\nolimits}
\def\proetale{{pro\text{-}\acute{e}tale}}
\def\etale{{\acute{e}tale}}
\def\QCoh{\mathit{QCoh}}
\def\Ker{\mathop{\mathrm{Ker}}}
\def\Im{\mathop{\mathrm{Im}}}
\def\Coker{\mathop{\mathrm{Coker}}}
\def\Coim{\mathop{\mathrm{Coim}}}

% Boxtimes
%
\DeclareMathSymbol{\boxtimes}{\mathbin}{AMSa}{"02}

%
% Macros for moduli stacks/spaces
%
\def\QCohstack{\mathcal{QC}\!\mathit{oh}}
\def\Cohstack{\mathcal{C}\!\mathit{oh}}
\def\Spacesstack{\mathcal{S}\!\mathit{paces}}
\def\Quotfunctor{\mathrm{Quot}}
\def\Hilbfunctor{\mathrm{Hilb}}
\def\Curvesstack{\mathcal{C}\!\mathit{urves}}
\def\Polarizedstack{\mathcal{P}\!\mathit{olarized}}
\def\Complexesstack{\mathcal{C}\!\mathit{omplexes}}
% \Pic is the operator that assigns to X its picard group, usage \Pic(X)
% \Picardstack_{X/B} denotes the Picard stack of X over B
% \Picardfunctor_{X/B} denotes the Picard functor of X over B
\def\Pic{\mathop{\mathrm{Pic}}\nolimits}
\def\Picardstack{\mathcal{P}\!\mathit{ic}}
\def\Picardfunctor{\mathrm{Pic}}
\def\Deformationcategory{\mathcal{D}\!\mathit{ef}}
\usepackage{cleveref}
\crefname{originalreference}{original reference}{original references}
\Crefname{originalreference}{Original reference}{Original references}
\makeatletter
\newlabel{situation-setup}{{02QL}{source}{}{sourcepointer-1}{}}
\newlabel{situation-setup@cref}{{[originalreference][0][]02QL}{[1][0][]source}{}{}{}}
\makeatother
\crefname{originallemma}{original lemma}{original lemmas}
\Crefname{originallemma}{Original lemma}{Original lemmas}
\makeatletter
\newlabel{lemma-localized-chern-pre}{{0F9K}{source}{}{sourcepointer-2}{}}
\newlabel{lemma-localized-chern-pre@cref}{{[originallemma][0][]0F9K}{[1][0][]source}{}{}{}}
\makeatother
\crefname{originalreference}{original reference}{original references}
\Crefname{originalreference}{Original reference}{Original references}
\makeatletter
\newlabel{situation-loc-chern}{{0GUJ}{source}{}{sourcepointer-3}{}}
\newlabel{situation-loc-chern@cref}{{[originalreference][0][]0GUJ}{[1][0][]source}{}{}{}}
\makeatother
\crefname{originallemma}{original lemma}{original lemmas}
\Crefname{originallemma}{Original lemma}{Original lemmas}
\makeatletter
\newlabel{lemma-base-change-localized-chern-pre}{{0GUI}{source}{}{sourcepointer-4}{}}
\newlabel{lemma-base-change-localized-chern-pre@cref}{{[originallemma][0][]0GUI}{[1][0][]source}{}{}{}}
\makeatother
\crefname{originallemma}{original lemma}{original lemmas}
\Crefname{originallemma}{Original lemma}{Original lemmas}
\makeatletter
\newlabel{lemma-gysin-at-infty}{{0F9J}{source}{}{sourcepointer-5}{}}
\newlabel{lemma-gysin-at-infty@cref}{{[originallemma][0][]0F9J}{[1][0][]source}{}{}{}}
\makeatother
\crefname{originallemma}{original lemma}{original lemmas}
\Crefname{originallemma}{Original lemma}{Original lemmas}
\makeatletter
\newlabel{lemma-base-change-silly}{{0GUG}{source}{}{sourcepointer-6}{}}
\newlabel{lemma-base-change-silly@cref}{{[originallemma][0][]0GUG}{[1][0][]source}{}{}{}}
\makeatother
\crefname{originallemma}{original lemma}{original lemmas}
\Crefname{originallemma}{Original lemma}{Original lemmas}
\makeatletter
\newlabel{lemma-homomorphism-pre}{{0FAR}{source}{}{sourcepointer-7}{}}
\newlabel{lemma-homomorphism-pre@cref}{{[originallemma][0][]0FAR}{[1][0][]source}{}{}{}}
\makeatother
\crefname{originallemma}{original lemma}{original lemmas}
\Crefname{originallemma}{Original lemma}{Original lemmas}
\makeatletter
\newlabel{lemma-loc-chern-agree}{{0FB8}{source}{}{sourcepointer-8}{}}
\newlabel{lemma-loc-chern-agree@cref}{{[originallemma][0][]0FB8}{[1][0][]source}{}{}{}}
\makeatother
\title{A composite of localized Chern classes equals the pushforward of the restricted classes along the intersection of the W\_i}
\author{The Stacks Project Authors}
\date{stacks-project:chow.tex:local-snapshot-20261012; GNU FDL 1.2 or later}
\begin{document}
\maketitle
\noindent\textit{Editorial note.} All delivered mathematical text is the author's own, copied from the pinned Stacks Project chapter file chow.tex; no statement or proof step was rewritten, reordered or removed.\par\medskip
\noindent\textit{Editorial note.} Preamble provenance: the pinned chapter file opens with an input of the shared preamble and carries no preamble text of its own, so the project's single Stacks-wide preamble staged at sources/stacks-chow/preamble.tex is excerpted verbatim as two spans, lines 14-18 (Xy-pic setup) and lines 154-217 (the native theorem declarations and mathematical macros). Nothing was added, removed or redefined; the wrapper class and the amsmath/amssymb/amsthm/graphicx base packages are supplied by the assembler.\par\medskip
\noindent\textit{Editorial note.} The counted statement and proof contain no cross reference, no citation and no included file, so the delivered body is self contained.\par\medskip
\noindent\textit{Editorial note.} Canonical location of the counted proposition: Stacks tag 0FE7 (sources/stacks-chow/pinned-original-source.tex line 11026).\par\medskip
\noindent\textit{Allowed standard prerequisites (editorial).}\begin{itemize}
\item Localized Chern classes c'\_p(Q) and the operators P'\_p(Q) attached to a perfect object and a closed subset. Same chapter, Lemma labelled lemma-localized-chern-pre, invoked in the statement and cited at proof lines 11062-11067. Supplies the objects being composed and the case n = 1 identification.
\item Base change behaviour of the localized class under restriction to a closed subscheme. Same chapter, Lemma labelled lemma-base-change-localized-chern-pre, cited at proof lines 11078-11082. Computes the restriction of c'\_\{p\_n\}(Q\_n) to Z\_\{n-1\} cap ... cap Z\_1.
\item The class C of the inclusion at infinity and its behaviour with pushforward. Same chapter, Lemma labelled lemma-gysin-at-infinity, cited at proof lines 11090-11091, and Lemma labelled lemma-homomorphism-pre, cited at proof lines 11108-11113. Provides the class C\_\{n-1\} and the identity C\_\{n-1\} o pushforward = id on Gysin images that reduces the lemma to a statement about images.
\item Agreement of the restricted localized class with the bivariant Chern class of W\_i -> W. Same chapter, Lemma labelled lemma-loc-chern-agree, cited at proof lines 11130-11134. Writes the restricted classes as restrictions of bivariant classes, so that commutation with i\_infinity\textasciicircum{}* applies.
\end{itemize}\medskip
\ifdefined\hypertarget\hypertarget{sourcepointer-1}{}\fi
\noindent\textit{Source reference.} Original reference 02QL: sources/stacks-chow/pinned-original-source.tex, source lines 1372--1372. Same-chapter pointer into the author environment situation at source line 1372; invoked at proof line None of the counted proof. It is retained only as a navigation pointer; every load-bearing step is inside the delivered span.\par\medskip
\ifdefined\hypertarget\hypertarget{sourcepointer-2}{}\fi
\noindent\textit{Source reference.} Original lemma 0F9K: sources/stacks-chow/pinned-original-source.tex, source lines 10204--10204. Same-chapter pointer into the author environment lemma at source line 10204; invoked at proof line None of the counted proof. It is retained only as a navigation pointer; every load-bearing step is inside the delivered span.\par\medskip
\ifdefined\hypertarget\hypertarget{sourcepointer-3}{}\fi
\noindent\textit{Source reference.} Original reference 0GUJ: sources/stacks-chow/pinned-original-source.tex, source lines 10678--10678. Same-chapter pointer into the author environment situation at source line 10678; invoked at proof line None of the counted proof. It is retained only as a navigation pointer; every load-bearing step is inside the delivered span.\par\medskip
\ifdefined\hypertarget\hypertarget{sourcepointer-4}{}\fi
\noindent\textit{Source reference.} Original lemma 0GUI: sources/stacks-chow/pinned-original-source.tex, source lines 10276--10276. Same-chapter pointer into the author environment lemma at source line 10276; invoked at proof line 11076 of the counted proof. It is retained only as a navigation pointer; every load-bearing step is inside the delivered span.\par\medskip
\ifdefined\hypertarget\hypertarget{sourcepointer-5}{}\fi
\noindent\textit{Source reference.} Original lemma 0F9J: sources/stacks-chow/pinned-original-source.tex, source lines 10026--10026. Same-chapter pointer into the author environment lemma at source line 10026; invoked at proof line 11086 of the counted proof. It is retained only as a navigation pointer; every load-bearing step is inside the delivered span.\par\medskip
\ifdefined\hypertarget\hypertarget{sourcepointer-6}{}\fi
\noindent\textit{Source reference.} Original lemma 0GUG: sources/stacks-chow/pinned-original-source.tex, source lines 9775--9775. Same-chapter pointer into the author environment lemma at source line 9775; invoked at proof line 11103 of the counted proof. It is retained only as a navigation pointer; every load-bearing step is inside the delivered span.\par\medskip
\ifdefined\hypertarget\hypertarget{sourcepointer-7}{}\fi
\noindent\textit{Source reference.} Original lemma 0FAR: sources/stacks-chow/pinned-original-source.tex, source lines 10132--10132. Same-chapter pointer into the author environment lemma at source line 10132; invoked at proof line 11117 of the counted proof. It is retained only as a navigation pointer; every load-bearing step is inside the delivered span.\par\medskip
\ifdefined\hypertarget\hypertarget{sourcepointer-8}{}\fi
\noindent\textit{Source reference.} Original lemma 0FB8: sources/stacks-chow/pinned-original-source.tex, source lines 10975--10975. Same-chapter pointer into the author environment lemma at source line 10975; invoked at proof line 11134 of the counted proof. It is retained only as a navigation pointer; every load-bearing step is inside the delivered span.\par\medskip
\begin{lemma}
\label{lemma-homomorphism-final}
Let $(S, \delta)$ be as in Situation \ref{situation-setup}. Let $X$ be
locally of finite type over $S$. Let $b : W \longrightarrow \mathbf{P}^1_X$
be a proper morphism of schemes. Let $n \geq 1$. For $i = 1, \ldots, n$
let $Z_i \subset X$ be a closed subscheme, let $Q_i \in D(\mathcal{O}_W)$
be a perfect object, let $p_i \geq 0$ be an integer, and let
$T_i \subset W_\infty$, $i = 1, \ldots, n$ be closed.
Denote $W_i = b^{-1}(\mathbf{P}^1_{Z_i})$. Assume
\begin{enumerate}
\item for $i = 1, \ldots, n$ the assumption of
Lemma \ref{lemma-localized-chern-pre} hold for
$b, Z_i, Q_i, T_i, p_i$,
\item $Q_i|_{W \setminus W_i}$ is zero, resp.\ isomorphic to a finite
locally free module of rank $< p_i$ placed in cohomological degree $0$,
\item $Q_i$ on $W$ satisfies
assumption (3) of Situation \ref{situation-loc-chern}.
\end{enumerate}
Then $P'_{p_n}(Q_n) \circ \ldots \circ P'_{p_1}(Q_1)$ is equal to
$$
(W_{n, \infty} \cap \ldots \cap W_{1, \infty} \to
Z_n \cap \ldots \cap Z_1)_* \circ
P'_{p_n}(Q_n|_{W_{n, \infty}}) \circ \ldots \circ P'_{p_1}(Q_1|_{W_{1, \infty}})
\circ C
$$
in $A^{p_n + \ldots + p_1}(Z_n \cap \ldots \cap Z_1 \to X)$,
resp.\ $c'_{p_n}(Q_n) \circ \ldots \circ c'_{p_1}(Q_1)$ is equal to
$$
(W_{n, \infty} \cap \ldots \cap W_{1, \infty} \to
Z_n \cap \ldots \cap Z_1)_* \circ
c'_{p_n}(Q_n|_{W_{n, \infty}}) \circ \ldots \circ c'_{p_1}(Q_1|_{W_{1, \infty}})
\circ C
$$
in $A^{p_n + \ldots + p_1}(Z_n \cap \ldots \cap Z_1 \to X)$.
\end{lemma}

\begin{proof}
Let us prove the statement on Chern classes by induction on $n$;
the statement on $P_p(-)$ is proved in the exact same manner.
The case $n = 1$ is the construction of $c'_{p_1}(Q_1)$ because
$W_{1, \infty}$ is the inverse image of $Z_1$ in $W_\infty$.
For $n > 1$ we have by induction that
$c'_{p_n}(Q_n) \circ \ldots \circ c'_{p_1}(Q_1)$ is equal to
$$
c'_{p_n}(Q_n) \circ
(W_{n - 1, \infty} \cap \ldots \cap W_{1, \infty} \to
Z_{n - 1} \cap \ldots \cap Z_1)_* \circ
c'_{p_{n - 1}}(Q_{n - 1}|_{W_{n - 1}, \infty}) \circ \ldots \circ
c'_{p_1}(Q_1|_{W_{1, \infty}})
\circ C
$$
By Lemma \ref{lemma-base-change-localized-chern-pre} the restriction of
$c'_{p_n}(Q_n)$ to $Z_{n - 1} \cap \ldots \cap Z_1$ is computed by
the closed subset $Z_n \cap \ldots \cap Z_1$, the morphism
$b' : W_{n - 1} \cap \ldots \cap W_1 \to
\mathbf{P}^1_{Z_{n - 1} \cap \ldots \cap Z_1}$
and the restriction of $Q_n$ to $W_{n - 1} \cap \ldots \cap W_1$.
Observe that $(b')^{-1}(Z_n) = W_n \cap \ldots \cap W_1$
and that $(W_n \cap \ldots \cap W_1)_\infty =
W_{n, \infty} \cap \ldots \cap W_{1, \infty}$.
Denote $C_{n - 1} \in A^0(W_{n - 1, \infty} \cap \ldots \cap W_{1, \infty} \to
Z_{n - 1} \cap \ldots \cap Z_1)$ the class of Lemma \ref{lemma-gysin-at-infty}.
We conclude the restriction of $c'_{p_n}(Q_n)$ to
$Z_{n - 1} \cap \ldots \cap Z_1$ is
\begin{align*}
&
(W_{n, \infty} \cap \ldots \cap W_{1, \infty} \to Z_n \cap \ldots \cap Z_1)_*
\circ
c'_{p_n}(Q_n|_{(W_n \cap \ldots \cap W_1)_\infty})
\circ
C_{n - 1} \\
& =
(W_{n, \infty} \cap \ldots \cap W_{1, \infty} \to Z_n \cap \ldots \cap Z_1)_*
\circ
c'_{p_n}(Q_n|_{W_{n, \infty}})
\circ
C_{n - 1}
\end{align*}
where the equality follows from Lemma \ref{lemma-base-change-silly}
(we omit writing the restriction on the right). Hence the above becomes
\begin{align*}
(W_{n, \infty} \cap \ldots \cap W_{1, \infty} \to
Z_n \cap \ldots \cap Z_1)_* \circ
c'_{p_n}(Q_n|_{W_n, \infty}) \circ \\
C_{n - 1} \circ
(W_{n - 1, \infty} \cap \ldots \cap W_{1, \infty} \to
Z_{n - 1} \cap \ldots \cap Z_1)_* \\
\circ
c'_{p_{n - 1}}(Q_{n - 1}|_{W_{n - 1}, \infty}) \circ \ldots \circ
c'_{p_1}(Q_1|_{W_{1, \infty}})
\circ C
\end{align*}
By Lemma \ref{lemma-homomorphism-pre}
we know that the composition
$C_{n - 1} \circ (W_{n - 1, \infty} \cap \ldots \cap W_{1, \infty} \to
Z_{n - 1} \cap \ldots \cap Z_1)_*$
is the identity on elements in the image of the Gysin map
$$
(W_{n - 1, \infty} \cap \ldots \cap W_{1, \infty} \to
W_{n - 1} \cap \ldots \cap W_1)^*
$$
Thus it suffices to show that any element in the image of
$c'_{p_{n - 1}}(Q_{n - 1}|_{W_{n - 1}, \infty}) \circ \ldots \circ
c'_{p_1}(Q_1|_{W_{1, \infty}}) \circ C$
is in the image of the Gysin map. We may write
$$
c'_{p_i}(Q_i|_{W_{i, \infty}}) = \text{restriction of } c_{p_i}(W_i \to W, Q_i)
\text{ to } W_{i, \infty}
$$
by Lemma \ref{lemma-loc-chern-agree} and assumptions (2) and (3) on $Q_i$
in the statement of the lemma. Thus, if $\beta \in \CH_{k + 1}(W)$
restricts to the flat pullback of $\alpha$ on $b^{-1}(\mathbf{A}^1_X)$,
then
\begin{align*}
& c'_{p_{n - 1}}(Q_{n - 1}|_{W_{n - 1}, \infty}) \cap \ldots \cap
c'_{p_1}(Q_1|_{W_{1, \infty}})
\cap C \cap \alpha \\
& =
c'_{p_{n - 1}}(Q_{n - 1}|_{W_{n - 1}, \infty}) \cap \ldots \cap
c'_{p_1}(Q_1|_{W_{1, \infty}})
\cap i_\infty^* \beta \\
& =
c_{p_{n - 1}}(W_{n - 1} \to W, Q_{n - 1}) \cap \ldots \cap
c_{p_{n - 1}}(W_1 \to W, Q_1) \cap i_\infty^* \beta \\
& =
(W_{n - 1, \infty} \cap \ldots \cap W_{1, \infty} \to
W_{n - 1} \cap \ldots \cap W_1)^*
\left(c_{p_{n - 1}}(W_{n - 1} \to W, Q_{n - 1}) \cap \ldots \cap
c_{p_1}(W_1 \to W, Q_1) \cap \beta\right)
\end{align*}
as desired. Namely, for the last equality we use that
$c_{p_i}(W_i \to W, Q_i)$ is a bivariant class and hence
commutes with $i_\infty^*$ by definition.
\end{proof}

\end{document}
